
où, si @@M0@@ est un
espace vectoriel sur @@M1@@, on écrit @@M2@@ pour @@M3@@. Pour @@M4@@, @@M5@@ est surjectif par (ii); donc, si @@M6@@ est surjectif, alors @@M7@@ l'est. Par le théorème de Serre, @@M8@@ est surjective pour @@M9@@; d'où (iii).

\textbf{Proposition 1.4.} \emph{Soit @@M10@@ une suite exacte sur @@M11@@. Si @@M12@@ est @@M13@@-régulier, alors on a}~:
\begin{enumerate}[label=(\roman*)]
\item \emph{@@M14@@ pour @@M15@@ et @@M16@@.}
\item \emph{@@M17@@ pour @@M18@@.}
\item \emph{@@M19@@ pour @@M20@@.}
\end{enumerate}
\emph{En particulier, @@M21@@ est @@M22@@-régulier.}

En effet, dans la suite exacte
@@M23@@
les deux groupes extrêmes sont @@M24@@ pour @@M25@@ et @@M26@@ par 1.3 (i). Donc @@M27@@, mais, par le théorème de Serre, @@M28@@ pour @@M29@@; d'où (i).

Pour @@M30@@, considérons la suite exacte suivante :
@@M31@@
Elle implique (ii). Or, supposons @@M32@@ surjectif et considérons le diagramme commutatif suivant :
@@M33@@
